% Plantilla para documentos LaTeX para enunciados
% Por Pedro Pablo Aste Kompen - ppaste@uc.cl
% Licencia Creative Commons BY-NC-SA 3.0
% http://creativecommons.org/licenses/by-nc-sa/3.0/

\documentclass[12pt]{article}

% Margen de 1 pulgada por lado
\usepackage{fullpage}
% Incluye gráficas
\usepackage{graphicx}
% Packages para matemáticas, por la American Mathematical Society
\usepackage{amssymb}
\usepackage{amsmath}
% Desactivar hyphenation
\usepackage[none]{hyphenat}
% Saltar entre párrafos - sin sangrías
\usepackage{parskip}
% Español y UTF-8
\usepackage[spanish]{babel}
\usepackage[utf8]{inputenc}
% Links en el documento
\usepackage[hidelinks]{hyperref}
\usepackage[shortlabels]{enumitem}

\usepackage{algorithmic}
\usepackage[nothing]{algorithm}

\newcommand{\bcap}{\ \hat{\cap} \ }
\newcommand{\bcup}{\ \hat{\cup} \ }
\newcommand{\bagP}{\hat{\mathcal{P}}}
\newcommand{\la}{\leftarrow}


\newcommand{\N}{\mathbb{N}}
\newcommand{\Z}{\mathbb{Z}}
\newcommand{\Exp}[1]{\mathcal{E}_{#1}}
\newcommand{\List}[1]{\mathcal{L}_{#1}}
\newcommand{\EN}{\Exp{\N}}
\newcommand{\LN}{\List{\N}}
\newcommand{\lrinline}{\text{\textbf{ --- }}}
\newcommand{\lrmatrix}{\text{\textbf{---}}}

\newcommand{\comment}[1]{}
\newcommand{\lb}{\\~\\}
\newcommand{\eop}{_{\square}}
\newcommand{\hsig}{\hat{\sigma}}
\newcommand{\ra}{\rightarrow}
\newcommand{\lra}{\leftrightarrow}
\newcommand{\com}{\textquotedbl}
\newcommand{\htau}{\hat{\tau}}
\newcommand{\twopartdef}[6]
{
	\left\{
		\begin{array}{ll}
			#1 &  #2 \\
			#3 &  #4 \\
			#5 &  #6
		\end{array}
	\right.
}

\newcommand{\mtwopartdef}[8]
{
	\left\{
		\begin{array}{ll}
			#1 &  #2 \\
			#3 &  #4 \\
			#5 &  #6 \\
			#7 &  #8
		\end{array}
	\right.
}

% Comentar aligns
\newcommand{\step}[1]{%
  \text{\phantom{(#1)}} \tag{#1}
}

\usepackage[sol]{optional}

\begin{document}
\thispagestyle{empty}
% Membrete
% PUC-ING-DCC-IIC1103
\begin{minipage}{2.3cm}
\includegraphics[width=2cm]{logo.pdf}
\vspace{0.5cm} % Altura de la corona del logo, así el texto queda alineado verticalmente con el círculo del logo.
\end{minipage}
\begin{minipage}{\linewidth}
\textsc{\raggedright \footnotesize
Pontificia Universidad Católica de Chile \\
Departamento de Ciencia de la Computación \\
IIC1253 - Matemáticas Discretas \\}
\end{minipage}


% Titulo
\begin{center}
\vspace{0.5cm}
{\huge\bf Tarea 4}\\
\vspace{0.2cm}
6 de octubre de 2020\\
\vspace{0.2cm}
\footnotesize{2º semestre 2020 - Profesores G. Diéguez - F. Suárez}
\vspace{0.2cm}
\rule{\textwidth}{0.05mm}
\end{center}

% Cuerpo del documento ejemplo

\section*{Requisitos}
\begin{itemize}

  \item La tarea es individual. Los casos de copia serán sancionados con la
   reprobación del curso con nota $1.1$.

  \item \textbf{Entrega}: Hasta las 23:59:59 del 19 de octubre a través del buzón habilitado en el sitio del curso (Canvas).

  \begin{itemize}

    \item Esta tarea debe ser hecha completamente en \LaTeX. Tareas hechas a mano o
    en otro procesador de texto \textbf{no serán corregidas}.

    \item Debe usar el template \LaTeX\ publicado en la página del curso.

    \item Cada problema debe entregarse en un archivo independiente de las demas preguntas.

    \item Los archivos que debe entregar son un archivo \texttt{PDF} por cada pregunta
    con su solución con nombre \texttt{numalumno-P1.pdf} y \texttt{numalumno-P2.pdf} ,
    junto con un \texttt{zip} con nombre \texttt{numalumno.zip}, conteniendo los archivos \texttt{numalumno-P1.tex} y \texttt{numalumno-P2.tex} que compilan su tarea.
    Si su código hace referencia a otros archivos, debe incluirlos también.

  \end{itemize}

  \item El no cumplimiento de alguna de las reglas se penalizará con un descuento de 0.5 en la nota final (acumulables).

  \item No se aceptarán tareas atrasadas.

  \item Si tiene alguna duda, el foro de Canvas es el lugar oficial para realizarla.
\end{itemize}


\newpage
\section*{Problemas}

    \subsection*{Problema 1 - Relaciones de orden}
      %!TEX root =./tarea2.tex

\newcommand\myeq{\stackrel{\mathclap{\scriptsize\mbox{HI}}}{\equiv}}

\begin{enumerate}[a)]
    
    \item Ocupando la definición de suma vista en clases:
    
    \hspace{10mm} 1. $sum(m, 0) = m$
    
    \hspace{10mm} 2. $sum(m, \sigma(n)) =\sigma(sum(m, n))$
    
    Es trivial concluir que el 0 es elemento neutro.
    
    \vspace{5mm}
    
    \item Ocupando la definición de multiplicación:
    
    \hspace{10mm} 1. $mult(m, 0) = 0$
    
    \hspace{10mm} 2. $mult(m, \sigma(n)) =sum(m, mult(m, n))$
    
    notemos que 
    \begin{align*}
        \hspace{-20mm} mult(m, 1) &= mult(m, \sigma(0)) \\
                   &= sum(m, mult(m, 0)) \\
                   &= sum(m, 0) = m
    \end{align*}
    
    Por lo que $1$ es elemento neutro de la multiplicación
    
    
    \item Queremos demostrar que 

    \hspace{20mm} $mult(a, sum(m, n)) = sum(mult(a, m), mult(a, n))$
    
    para cualquier $a, m, n$.
    
    Por inducción estructural:
    
    \textbf{BI}: Para $n = 0$ tenemos que
    \begin{align*}
        mult(a, sum(m, 0)) &= mult(a, m) \\
                           &= sum(mult(a, m), 0))\\
                           &= sum(mult(a, m), mult(a, 0))
    \end{align*}
    
    por lo que el caso base se cumple.
    
    \newpage
    
    \textbf{HI}: Supongamos que se cumple que 
    \begin{align*}
        mult(a, sum(m, n)) = sum(mult(a, m), mult(a, n))
    \end{align*}
    
    \textbf{TI}: demostrar que tambien se cumple que
    \begin{align*}
        mult(a, sum(m, \sigma(n))) &= sum(mult(a, m), mult(a, \sigma(n)) \\
                                   &= sum(mult(a, m), sum(a, mult(a, n)))
    \end{align*}
    
    Partiremos por el lado izquierdo:
    
    \begin{align*}
        mult(a, sum(m, \sigma(n))) &= mult(a, \sigma(sum(m, n))) \\
                                   &= sum(a, mult(a, sum(m, n))) \\
                                   &\myeq sum(a, sum(mult(a, m), mult(a, n))) \\
                                   &= sum(mult(a, m), sum(a, mult(a, n))) \text{ por asociatividad} \\
    \end{align*}
    
    con lo que se demuestra la tesis de inducción, y la suma es distributiva con la multiplicación.
\end{enumerate}

     \subsection*{Problema 2 - Funciones y cardinalidad}
      %!TEX root =./tarea4.tex
\begin{enumerate}[a)]
    \item Si definimos $C \in S$ como el conjunto de máxima cardinalidad posible de números no consecutivos, este será de tamaño $n$. Está definido de tal manera que no podemos añadir otro elemento en $S \setminus C$ a $C$, pues dejaría de ser no consecutivo. Ahora intentamos armar el conjunto $X$ a partir de $C$, de manera que $X$ sea no consecutivo. Como $\vert X \vert = n + 1$, no podemos mapear cada elemento de $C$ a $X$, por el principio del palomar (pues no existe una función sobreyectiva desde $C$ a $X$), por lo tanto nos faltaría agregar un elemento al conjunto $X$ para que su cardinalidad sea $n+1$. Pero como ya no quedan elementos del conjunto $C$, debemos agregar elementos de $S \setminus C$, con lo que $X$ dejaría de ser no consecutivo, y tendría almenos un par de números consecutivos. 

   
\end{enumerate}


\end{document}
